\subsection{Détermination de l'angle utile du prisme}

\begin{table}[ht]
    \centering
    
    \begin{tabular}{|c|c|c|c|c|}
        \hline
        Essai & $\alpha_1$ ($^\circ$) & $\alpha_2$ ($^\circ$) & $|\alpha_1 - \alpha_2|$ ($^\circ$) & $A = \frac{|\alpha_1 - \alpha_2|}{2}$ ($^\circ$) \\
        \hline
        Face 1 & 229 & 109 & 120 & 60,0 \\
        Face 2 & 242 & 121 & 121 & 60,5 \\
        Face 3 & 235 & 114 & 121 & 60,5 \\
        \hline
        \multicolumn{4}{|r|}{\textbf{Moyenne $\bar{A}$ ($^\circ$)}} & \textbf{60,3} \\
        \hline
    \end{tabular}
    \label{tab:mesure_angle_A}
    \caption{Détermination expérimentale de l'angle $A$ du prisme par réflexion}
\end{table}

L'objectif de cette partie est de présenter sous forme d'un tableau nos résultats de manière brute afin de calculer l'angle utile du prisme $A$.

Puisque \(L2 - L1 =  2A \) alors il apparait que :
\[A = \frac{\left | L1 - L2 \right |}{2} \]  


Compte tenu de la précision des lectures sur le cercle gradué ($1^{\circ}$), 
nous retenons une incertitude $\Delta A = 0.5^{\circ}$ donnant un résultat final de 
\[A = (60.3 \pm 0.5)^{\circ}\]

\subsection{Etude de la déviation \texorpdfstring{$D = f(\alpha_i)$}{D = f(alpha)}}

\begin{table}[h]
\centering
\begin{tabular}{|l|r|r|r|r|r|r|r|}
\hline
$\alpha_i$ & 50  & 55  & 60  & 65  & 70  & 75  & 80  \\ \hline
$L$        & 102 & 112 & 122 & 132 & 142 & 152 & 162 \\ \hline
$\alpha_d$ & 239 & 237 & 236 & 238 & 241 & 245 & 247 \\ \hline
$D$        & 61  & 59  & 58  & 60  & 63  & 67  & 69  \\ \hline
\end{tabular}
\caption{Relevé de la caractéristique $D = f(\alpha_i)$. Les mesures sont relevées en degrés ($^{\circ}$)}
\end{table}
\begin{figure}[h]
    \centering
    \includestandalone{sections/resultats/plot_manip_b} % Sans l'extension .tex
    \caption{Évolution de la déviation en fonction de l'angle d'incidence.}
\end{figure}

Nous observons une courbe que l'on peut imaginer comme étant modélisable par une parabole. 
Etant donné la granularité dans nos mesures, notre courbe est plus "saccadée". 
Nous arrivons toutefois à déterminer par lecture graphique que le minimum de déviation 
$D_m = 58^{\circ}$ 

\subsection{Mesure précise du minimum de déviation \texorpdfstring{$D_m$}{Dm}}

\subsection{Calcul de l'indice de réfraction}